\chapter{SVAR com Restrições de Longo Prazo}
\label{cap:svar2}
% ── índice ──────────────────────────────────────────────────────
\index{Blanchard-Quah, decomposição de|textbf}
\index{choque permanente|textbf}
\index{choque transitório|textbf}
\index{restrições de longo prazo|textbf}
\index{SVEC}

% ─────────────────────────────────────────────────────────────────────────────
\section{Identificação de Blanchard-Quah}

O cap.~\ref{cap:svar} (SVAR Cholesky) impõe restrições de \emph{curto prazo}:
a variável na posição 1 não é afetada contemporaneamente pela segunda.
Blanchard e Quah (1989) propõem uma alternativa via restrições de
\emph{longo prazo}: choques permanentes e transitórios.

A ideia central: na macroeconomia, um choque de oferta tem efeito permanente
sobre o produto ($y$), enquanto um choque de demanda tem efeito apenas
transitório. A restrição:
\[
  \sum_{h=0}^\infty \Theta_h^{(1,1)} = \text{livre}, \quad
  \sum_{h=0}^\infty \Theta_h^{(1,2)} = 0
\]
impõe que o choque 2 não afeta $y$ no longo prazo. Essa restrição identifica
a decomposição $B$ sem nenhuma restrição de impacto contemporâneo.

\subsection*{Decomposição BQ}

Dado o VAR na forma reduzida $y_t = \Phi_1 y_{t-1} + u_t$:
\begin{enumerate}
  \item Compute a soma $C = (I - \Phi_1)^{-1}$ — multiplicador de longo prazo
  \item Fatore $C\,\Sigma_u\,C' = F\,F'$ (Cholesky de $F$)
  \item Extraia $B = C^{-1} F$ — os choques estruturais de longo prazo
\end{enumerate}

% ─────────────────────────────────────────────────────────────────────────────
\section{Verificação por DGP}

DGP: VAR(1) bivariado onde $y_1$ é um passeio aleatório (choque 1 permanente),
$y_2$ é estacionária e responde transitoriamente ao choque 1.

\begin{lstlisting}[caption={SVAR Blanchard-Quah}]
seed(42)
// y1_t = y1_{t-1} + e1_t               (passeio aleatório)
// y2_t = 0.5*y2_{t-1} + 0.4*e1_t + 0.6*e2_t  (estacionária)

let m_bq = svar(df, y1, y2, lags=1, id=longrun)
display m_bq
let m_irf_bq = sirf(m_bq, steps=20)
display m_irf_bq
\end{lstlisting}

\subsection*{Matriz B (choques estruturais)}

\begin{lstlisting}[language={}, basicstyle=\ttfamily\small, frame=none,
  numbers=none, backgroundcolor=\color{codbg}]
SVAR(1) — Long-run restrictions
B Matrix:
[ 0.4907   0.0358 ]
[ 0.1935   0.3031 ]
\end{lstlisting}

A matriz $B$ não é triangular — ao contrário do Cholesky, as restrições de
longo prazo não implicam restrições de impacto contemporâneo.

\subsection*{IRF estrutural}

\begin{lstlisting}[language={}, basicstyle=\ttfamily\small, frame=none,
  numbers=none, backgroundcolor=\color{codbg}]
Impulso: Choque 1 (permanente)
  h= 0:  y1=+0.491  y2=+0.194
  h= 1:  y1=+0.464  y2=+0.087
  h= 5:  y1=+0.405  y2=-0.003   (efeito em y2 decai)
  h=19:  y1=+0.271  y2=-0.005   (y1 retém efeito permanente)

Impulso: Choque 2 (transitório)
  h= 0:  y1=+0.036  y2=+0.303
  h= 1:  y1=+0.016  y2=+0.144
  h= 5:  y1=-0.001  y2=+0.007   (efeito em y1 decai a ~0)
  h=19:  y1=-0.001  y2= 0.000   (transitório confirmado) ✓
\end{lstlisting}

O choque 2 afeta $y_1$ apenas transitoriamente (soma → 0), confirmando a
restrição de longo prazo. O choque 1 tem efeito acumulado crescente sobre $y_1$
— coerente com a raiz unitária.

% ─────────────────────────────────────────────────────────────────────────────
\section{Sintaxe}

\begin{lstlisting}
svar(df, y1, y2, lags=1, id=longrun)     // BQ: restrições de longo prazo
svar(df, y1, y2, lags=1, id=cholesky)    // Cholesky: restrições contemporâneas
sirf(m_svar, steps=20)                   // IRF estrutural
\end{lstlisting}

\begin{nota}
  A identificação BQ requer que as variáveis sejam integradas de ordem 1
  ($I(1)$). Para séries estacionárias, $(I - \hat{\Phi}_1)^{-1}$ é finito e
  a decomposição permanente/transitória não tem interpretação natural.
  Verificar raiz unitária (cap.~\ref{cap:arima}, \ref{cap:raiz_unitaria2})
  antes de aplicar BQ.
\end{nota}
